\documentclass[11pt,a4paper]{article}
\usepackage[T1]{fontenc}
\usepackage[utf8]{inputenc}
\usepackage{amsmath,amssymb,amsthm}
\usepackage[margin=1in]{geometry}
\usepackage[colorlinks=true,linkcolor=blue,urlcolor=blue]{hyperref}
\theoremstyle{plain}
\newtheorem{theorem}{Theorem}
\newtheorem{lemma}[theorem]{Lemma}
\newtheorem{proposition}[theorem]{Proposition}
\newtheorem{corollary}[theorem]{Corollary}
\theoremstyle{definition}
\newtheorem{remark}[theorem]{Remark}

\title{Recovering the unwritten recurrences of the table OEIS A236884}
\author{Adrian Perez Fontelles\\ \small Independent researcher}
\date{2 September 2026}
\begin{document}
\maketitle

\begin{abstract}
OEIS A236884 is a two-dimensional table $T(n,k)$ whose empirical recurrences are, for several of
its lines, given only as an ORDER --- ``k=2: [order 31]'' and the like --- with the
coefficients in a linked file rather than in the entry. Those conjectures can still be settled,
and the recurrences recovered. A line of the table is a fixed-width or fixed-height array
count, hence a walk count on finitely many vertices, hence satisfies some monic recurrence of
order at most that number; Berlekamp--Massey on twice as many exact terms returns the MINIMAL
one. Where its order equals the order the entry states --- and where the same holds after the
first few terms are dropped --- any recurrence of that order the line satisfies has a
characteristic polynomial that is a multiple of the minimal one and of equal degree, hence is
that polynomial. This note settles column 2, column 3.
\end{abstract}

\noindent\small 2020 Mathematics Subject Classification. 05A15, 11B37, 15A18.\normalsize

\section{The table and the unwritten conjectures}

OEIS A236884 is ``T(n,k) = Number of (n+1) X (k+1) 0..2 arrays with the maximum plus the upper median plus the lower median of every 2 X 2 subblock equal.''. It is read by antidiagonals and begins
\[
81, 211, 211, 537, 669, 537, 1523, 2111,\ \dots
\]
The entry carries, among its empirical claims:
\begin{quote}\small
k=2: [order 31]\\
k=3: [order 63]
\end{quote}
As of the ``Last modified'' line on the live entry (Jun 18 2022 18:00:45, revision 6) these are still
recorded as empirical, and the recurrences themselves are not in the entry text.

\section{Why a line of the table is a walk count}

Fix the index that the line fixes. The entry defines $T(n,k)$ as a number of arrays of a shape
determined by $n$ and $k$, subject to a condition local to a bounded window of consecutive
lines. Substituting the fixed index into the entry's own wording turns the definition into an
ordinary one-parameter array count, and for such a count the admissible windows are the
vertices of a finite digraph and an array is a walk. So the line is a walk count on $S$
vertices, and by the Cayley--Hamilton theorem it satisfies a monic integer recurrence of order
at most $S$.

\section{Recovering the recurrences}

\begin{lemma}\label{lem:bm}
A sequence known to satisfy some linear recurrence of order at most $S$ is determined, with its
minimal recurrence, by its first $2S$ terms; Berlekamp--Massey applied to them returns that
minimal recurrence exactly.
\end{lemma}

\subsection*{Column $k=2$}
Substituting the fixed index into the entry's wording gives an ordinary one-parameter array
count, a walk on $S=102$ vertices. Berlekamp--Massey on $2S=204$ exact terms
returns a minimal recurrence of order $31$ --- the order the entry states --- with
integer coefficients
\[
(c_1,\dots,c_{31})=(14,\ -28,\ -459,\ 2258,\ 4418,\ -46337,\ 15855,\ 478245,\ -699780,\ -2810786,\ 6844631,\ 9089948,\ -36536014,\ -9754145,\ 120336673,\ -40121314,\ -248844918,\ 191784528,\ 308729988,\ -375062032,\ -194512520,\ 397632272,\ 14941536,\ -230210592,\ 50723072,\ 68301056,\ -26409216,\ -8429056,\ 4614144,\ 122880,\ -184320).
\]
so that $a(m)=\sum_{i=1}^{31}c_i\,a(m-i)$ for every $m$. Removing the first
$31+4$ terms and repeating gives order $31$ again, which is what the
identification below needs.

\subsection*{Column $k=3$}
Substituting the fixed index into the entry's wording gives an ordinary one-parameter array
count, a walk on $S=251$ vertices. Berlekamp--Massey on $2S=502$ exact terms
returns a minimal recurrence of order $63$ --- the order the entry states --- with
integer coefficients
\[
(c_1,\dots,c_{63})=(16,\ 27,\ -1667,\ 4072,\ 76934,\ -341187,\ -2049421,\ 13078549,\ 33627314,\ -318556323,\ -306610231,\ 5496429956,\ 54510153,\ -70935843916,\ 45538376655,\ 706944602204,\ -806626246425,\ -5549283296459,\ 8679696177536,\ 34723994107362,\ -68113932764134,\ -174240406972383,\ 414218469332205,\ 701207183888323,\ -2008642190732129,\ -2246830446293595,\ 7889578094693248,\ 5619591657864136,\ -25328063997365173,\ -10441277176211870,\ 66791970202879633,\ 12296932132110369,\ -145011176324599384,\ -1055842701605466,\ 259241512822452234,\ -34025039397003762,\ -381021968390179388,\ 91434397011877932,\ 458967239587347048,\ -148931424202336496,\ -451028233509798832,\ 175195326671834064,\ 359397679149116544,\ -155964021652045664,\ -230443907502100160,\ 106480622832011904,\ 117774941203385216,\ -55687378365970432,\ -47417347574292480,\ 22079229293202432,\ 14817525871303680,\ -6514237146490880,\ -3524453417234432,\ 1390125346611200,\ 620930322202624,\ -205596162195456,\ -77774770995200,\ 19696234659840,\ 6478165180416,\ -1082279854080,\ -317278126080,\ 25480396800,\ 6794772480).
\]
so that $a(m)=\sum_{i=1}^{63}c_i\,a(m-i)$ for every $m$. Removing the first
$63+4$ terms and repeating gives order $63$ again, which is what the
identification below needs.

\begin{theorem}
For each line above, the displayed recurrence holds for every index, and it is the recurrence
the entry records.
\end{theorem}

\begin{proof}
It holds by Lemma~\ref{lem:bm}. For the identification, the entry asserts that the line
satisfies SOME recurrence of the stated order, possibly only from some index on. Let $q$ be its
characteristic polynomial and $q_{\min}$ the minimal polynomial of the tail. Then
$q_{\min}\mid q$, and the second Berlekamp--Massey run --- on the line with its first few
terms removed --- shows $\deg q_{\min}$ equals the stated order, which is $\deg q$. Both are
monic, so $q=q_{\min}$.
\end{proof}

\section{Verification}

Three checks.

First, each line's model was compared against the entry's own published values for that line,
taken from the antidiagonal DATA read in either orientation and from the ``Table starts''
block, and agrees at every available term. This is the only point at which reading the entry's
English enters, and a misreading gives different counts.

Second, each recovered recurrence was evaluated directly on the model's values, with no
Berlekamp--Massey involved, and holds at every index tested.

Third, each order was computed twice by independent routes --- modulo a $61$-bit prime, where
every intermediate is a machine word, and again in exact rational arithmetic --- and once more
on the line with its first few terms dropped, which is what the identification requires. All
agree.

\begin{thebibliography}{9}
\bibitem{oeis} The OEIS Foundation, \emph{The On-Line Encyclopedia of Integer
Sequences}, \url{https://oeis.org}, 2026. Sequence A236884.
\bibitem{massey} J.~L.~Massey, Shift-register synthesis and BCH decoding, \emph{IEEE
Trans. Inform. Theory} \textbf{15} (1969), 122--127.
\bibitem{stanley} R.~P.~Stanley, \emph{Enumerative Combinatorics}, Volume 1, 2nd ed.,
Cambridge University Press, 2011. (Transfer-matrix method, Section 4.7.)
\bibitem{horn} R.~A.~Horn and C.~R.~Johnson, \emph{Matrix Analysis}, 2nd ed., Cambridge
University Press, 2013.
\end{thebibliography}

\end{document}
